\documentclass[noauthor, nooutcomes]{ximera}

\input{../preamble}


\title{Fractions and Division}
\author{Vic Ferdinand, Betsy McNeal, Jenny Sheldon}

\begin{document}
\begin{abstract} 
In this activity, we try to understand how fractions and division are related.  Can all fractions be written as decimals?  Can all decimals be written as fractions?
\end{abstract}
\maketitle



\begin{problem}
Write a how many in one group story problem for $1 \div 6$, and explain how you know your story problem is correct. Then, use our meaning of fractions to explain why the answer to your story problem is $\frac{1}{6}$.

\end{problem}

\begin{problem}
Why does the previous problem mean that we can use long division to find the decimal representation for $\frac{1}{6}$?
\end{problem}



\begin{problem}
%we get R=0
Use long division to find the decimal representation for $\frac{1}{25}$. What happened? Was this a special case or will this happen for all fractions with a terminating decimal representation? Why?
\end{problem}



\begin{problem}
%not a remainder of zero
Use long division to find the decimal representation for $\frac{1}{7}$. What type of pattern occurred?  Why did this happen?  Was this a special case or will this happen for all fractions with a repeating decimal representation?  Why?
\end{problem}



\begin{problem}
Use long division to find the decimal representation for $\frac{1}{9}$. What type of pattern occurred?  Why did this happen?  Was this a special case or will this happen for all fractions with a repeating decimal representation?  Why?
\end{problem}



\begin{problem}
What do the previous three problems tell us about what the decimal representation of a fraction has to look like? Explain how you know. %When it comes to fractions (more technically rational numbers), there are two types: terminating and repeating. Look up a definition for each category, and then use that definition to give several examples.
\end{problem}


\begin{problem}
Are there any decimals that can't be written as fractions? Use your work in the previous problem to explain how you know.
\end{problem}



\begin{problem}
Bonus problem! Can you use what you already know about the decimal for $\frac{1}{9}$ to write the decimal for $\frac{4}{9}$? What about $\frac{1}{18}$? 
\end{problem}





 \newpage
\begin{instructorNotes}

{\bf Main Goal:} We would like students to come away from this activity understanding why we can use long division to find the decimal equivalent of a fraction through the use of an example.

{\bf Overall Picture:} We are drawing together many ideas in the course in this activity; don't be afraid to spend time if students need review on any topic.

 \begin{itemize}
 
 \item Connecting long division could potentially be a review problem depending on our work with division, but it's good to discuss in detail here. The idea is to see that solving the how many in one group problem involves splitting up the single item into six equal pieces, and placing one of these pieces in each group. Once we understand this connection, the answer to a long division problem is the same as the fraction, even if we represent it as a decimal. Note that if the numerator were not 1 in this case, we would be in the situation of the pies problem.
 
 \item If you have students who need an extra challenge, you could ask them to repeat the problem with a how many groups perspective!
 
 \item When we use long division to find decimal representations, there are several ideas to keep in mind. The most important part of these two problems is that students can execute the long division and explain why this gives us the decimal associated with the fraction (using the previous arguments). If we have time, it's interesting to explore some more ideas here, such as how the terminating or repeating process interacts with the long division (remainder of zero versus never getting a remainder of zero). We can also explore the fact that the repeating decimal has a pattern of remainders. See the next bullet point for optional ideas to discuss if time.
 
  \item In the problem where we divide $1 \div 7$, there are three points that need to be made to complete the argument:  That the division process is necessarily infinite (no zero remainder because of the 7 in the denominator -- we have not justified this, but students may be able to see this pattern), the number of possible remainders is finite and thus we will necessarily run into a remainder we have seen before, and finally, because we are dropping a zero down each time (leaving a multiple of ten to divide into), the sequence of divisions will be the same, generating the same sequence of quotients. This argument is difficult for students to put together (most have some  pieces, some all of them, but making a full coherent argument is difficult).
  
  \item After we do several division problems, you want to use the question about finding a pattern to bring up the idea of terminating and repeating decimals. Students should be able to express this in terms of whether we ever get a remainder of zero or whether the pattern will continue forever. (Since we didn't do the division forever, you want to make sure they feel very convinced that this pattern will continue without end!) We want to introduce the terminating and repeating terminology here.


\item The penultimate question gives us a chance to talk briefly about irrational numbers based on the fact that we've just said that fractions either terminate or repeat. Depending on time, you should have the students come up with examples of irrationals. The most common first examples are famous irrationals like $\pi$, $e$, $\sqrt{2}$. You can also encourage students to come up with a ``homemade'' irrational number such as $0.01001000100001\dots$, which encourages students to think about patterns that are not repeating, and perhaps also helps students to understand that there are a lot more irrational numbers out there than they perhaps imagined! 


 \end{itemize}


{\bf Suggested Timing:} We have two class periods to work through these ideas. Start by giving students 5-10 minutes to think about the first two problems, then spend about 20 minutes in discussion with groups presenting their work. This is the main point of the activity. Next, give students about 10 minutes to think about finding the decimals,  then discuss as much of the rest of the activity as you can. 
\end{instructorNotes}





\end{document}
